\documentclass{article}
\usepackage{graphicx} % Required for inserting images
\usepackage[svgnames]{xcolor}
\usepackage[paperwidth=1000mm, paperheight=1200mm, margin=3cm]{geometry}
\usepackage{anyfontsize}
\usepackage{tikz}
\usepackage{mathpazo}
\usepackage{multicol}
\usepackage{blindtext}

\renewcommand{\section}[1]{
    \begin{tikzpicture}
            \draw node[fill=yellow!20, text width=0.9\linewidth,
                text centered, inner sep=30pt, rounded corners=5pt,
                draw=yellow!80]
                {
                        \centering
                        \textbf{#1}
                };
        \end{tikzpicture}
}

\columnsep=60pt
\columnseprule=2pt

\begin{document}
    \begin{center}
        \fontsize{65}{75}
        \selectfont
        \begin{tikzpicture}
            \draw node[fill=yellow!20, text width=0.95\linewidth,
                text centered, inner sep=30pt, rounded corners=5pt,
                draw=yellow!80]
                {
                    \begin{minipage}{0.20\textwidth}
                        \includegraphics[width=0.6\textwidth]{Flag-logo.png}
                    \end{minipage}%
                    \begin{minipage}{0.55\textwidth}
                        \centering
                        Development of lattice-based abstractions in functional programs\\[1cm]
                        \fontsize{50}{60}
                        \selectfont
                        \textbf{Daniel Barrero} \hspace{0.5cm} \textbf{Valérie Gauthier} \hspace{0.5cm} \textbf{Nicolás Cardozo}\\
                        \fontsize{40}{50}
                        \selectfont
                        Systems and Computing Engineering  - Universidad de los Andes\\
                        \texttt{dr.barrero2562@uniandes.edu.co}
                    \end{minipage}%
                    \begin{minipage}{0.25\textwidth}
                    \centering
                        \includegraphics[width=0.5\textwidth]{logo_uniandes.jpg}
                    \end{minipage}
                    
                };
        \end{tikzpicture}
    \end{center}
    \vspace{3cm}
    \begin{multicols*}{3}
        \fontsize{50}{60}
        \selectfont
        \section{Quantum-safe Cryptography}
        Shor's quantum algorithm, published in 1994, demonstrated that factorization-based cryptography is no longer safe in the scenario of widespread availability of quantum computers.

        \bigskip
        
        \[
    S_{i,j} =
        \begin{array}{ccccccc}
           \          & \      & |00\rangle & |01\rangle & |10\rangle & |11\rangle & \ \\
           |00\rangle & \vline & 1          & 0          & 0          & 0          & \vline\\
           |01\rangle & \vline & 0          & 1          & 0          & 0          & \vline\\
           |10\rangle & \vline & 0          & 0          & 1          & 0          & \vline\\
           |11\rangle & \vline & 0          & 0          & 0          & e^{i\theta_{k-j}}          & \vline
        \end{array}
\]
    There is an \textbf{urgent} need for quantum-safe encryption protocols
    
    
    
    \end{multicols*}    
        
\end{document}

\endinput

Poner esto al final
    \begin{flushright}
      \includegraphics[width=0.3\columnwidth]{pil-logo}
      \includegraphics[width=0.1\columnwidth]{crypto_0}
    \end{flushright}
